\documentclass[10pt,a4paper]{amsart}
\usepackage[margin=19mm]{geometry}
\usepackage{amsmath,amssymb,mathtools}
\usepackage[scr=rsfs,cal=euler]{mathalfa}
\usepackage{xfrac}
\usepackage[unicode,hidelinks,bookmarks=false]{hyperref}
\newcommand{\ie}{i.e.\ }
\newcommand{\cf}{cf.\ }
\newcommand{\resp}{respectively\ }
\newcommand{\Subsection}{\subsection}
\providecommand{\nobreakdash}{\nobreak\hbox{-}}
\setlength{\emergencystretch}{2em}
\multlinegap0pt
\usepackage{bm}
\usepackage{bbm}
\usepackage{dsfont}
\usepackage{xspace}
\usepackage{cancel}
\usepackage{mathtools}
\usepackage{amsthm}
\makeatletter
\@ifundefined{theorem}{\newtheorem{theorem}{Theorem}}{}
\@ifundefined{lemma}{\newtheorem{lemma}[theorem]{Lemma}}{}
\makeatother
\newcommand{\E}{\mathbb{E}}
\newcommand{\bigO}{\mathcal{O}}
\providecommand{\N}{}
\providecommand{\P}{}
\title[Direct Paths in the Temporal Hypercube]{Direct Paths in the Temporal Hypercube}
\author{Austin Eide, Martijn G\"osgens, Pawe{\l} Pra{\l}at}
\date{Source: arXiv:2509.18931v1 (CC BY 4.0)}
\begin{document}
\maketitle

\noindent\emph{Source and licence.} Direct Paths in the Temporal Hypercube. Version arXiv:2509.18931v1 (CC BY 4.0), distributed under the Creative Commons Attribution 4.0 International licence, as recorded in the arXiv licence metadata for this exact version and shown on the abstract page https://arxiv.org/abs/2509.18931v1. Full rights notice: licenses/arxiv-2509.18931-CC-BY-4.0.txt.\par\smallskip

\noindent\textit{Editorial scope.} Counted unit: the Lemma whose source label is \texttt{lem:exp-limit2} in the article (source0.tex lines 631--637, source label \texttt{lem:exp-limit2}), delivered together with the article's own complete original proof of it (source0.tex lines 639--723, one \texttt{proof} environment of 85 lines). No mathematics is written, repaired or re-derived: statement and proof are the article's own text line for line. Required context retained uncounted: none. Every definition, display and label of those lines is retained unchanged. Omitted from this package and not redistributed: 1--630, 638--638, 724--1351. Nothing inside a retained excerpt is omitted or altered: every retained body is delivered exactly as the article's lines are written, so no omission receipt of any kind accompanies this package. Citations: the retained excerpts carry no citation command at all, so no bibliography entry of the article is transcribed and no reference is redistributed. Reference adaptation: none was needed and none was made; every reference inside the delivered unit resolves inside it, so no unresolved reference or citation is produced. Printed slips kept literal: no printed word, symbol, hypothesis or number of the counted statement or of its proof was altered. Layout and class: the article's own class is \texttt{article} with packages including hyperref, ulem, amsthm, amsmath, amssymb, bm, mathrsfs, amsfonts, mathtools, graphicx, tikz, enumitem, caption, titlesec; the wrapper is the lane's pinned \texttt{amsart} editorial wrapper at 10pt on A4, which restates the theorem-like environments the retained text needs and adds the article's own macro definitions that the retained text needs (\texttt{\textbackslash E}, \texttt{\textbackslash N}, \texttt{\textbackslash P}, \texttt{\textbackslash bigO}), with extra wrapper packages (bm, bbm, dsfont, xspace, cancel, mathtools). The delivered numbering is the unit's own. Assets: no retained span contains a figure environment, a graphics-inclusion command, a PDF-page inclusion, a table or a TikZ picture; no image file is copied into this package, the package declares no asset dependency and no image file is redistributed with it. Licence: version arXiv:2509.18931v1 of this article is distributed under the Creative Commons Attribution 4.0 International licence, as recorded in the arXiv licence metadata for this exact version and shown on the abstract page https://arxiv.org/abs/2509.18931v1. A full rights notice is delivered at licenses/arxiv-2509.18931-CC-BY-4.0.txt. Characters: every retained excerpt is pure ASCII, so the delivered engine, XeLaTeX with its default Latin Modern text font, reports no missing character.
\begin{lemma}\label{lem:exp-limit2}
    There exists a random variable $Z\sim\textup{Exp}(1)$, so that
    \(
    Z_{\infty,k,r_k}\stackrel{\P}\to Z,
    \)
    as $k\to\infty$ and ${r_k=\lceil\frac32\log k\rceil}.$
\end{lemma}

\begin{proof}
For $i \in [r]$, define
\[
Z_{k-1,r}^{(i)}=\sum_{(i_2,\dots,i_k)\in[r]^{k-1}}\exp\left(-\sum_{\ell=2}^k\sum_{j=1}^{i_{\ell}}\Delta(i,i_2,\dots,i_{\ell-1},j)\right).
\]
Then note that $Z_{k-1,r}^{(1)},\dots, Z_{k-1,r}^{(r)}$ are independent copies of $Z_{\infty,k-1,r}$, since they are functions of disjoint sets of $\Delta(\bm i)$'s.
Moreover,
\[
Z_{\infty,k,r}= \sum_{i=1}^r e^{-\sum_{j=1}^i\Delta(j)}Z_{k-1,r}^{(i)}.
\]
This tells us that
\[
\E[Z_{\infty,k,r_k}]=(1-2^{-r_k})\E[Z_{\infty,k-1,r_k}]=(1-2^{-r_k})^k\to1.
\]

We now show that the limiting random variable
\(
Z=\lim_{k\to\infty}\lim_{r\to\infty}Z_{\infty,k,r}
\)
exists. 
To show that the inner limit exists, note that $[r]^k\subset [r+1]^k$, so that $Z_{\infty,k,r+1}-Z_{\infty,k,r}>0$. Moreover, the Markov inequality allows us to bound
\begin{eqnarray*}
\P\left(Z_{\infty,k,r+1}-Z_{\infty,k,r}>r^{-2}\right) &\le& r^2\E[Z_{\infty,k,r+m}-Z_{\infty,k,r}] \\
&=& r^2\left((1-2^{-r-1})^k-(1-2^{-r})^k\right)=\bigO(r^2k2^{-r}).
\end{eqnarray*}
This yields
\[
\sum_{r=1}^\infty\P\left(Z_{\infty,k,r+1}-Z_{\infty,k,r}>r^{-2}\right) = k\sum_{r=1}^\infty \bigO(r^22^{-r}) <\infty.
\]
We then use the Borel-Cantelli lemma to conclude that $Z_{\infty,k,r}$ converges almost surely as $r\to\infty$ to some random variable $Z_{\infty,k,\infty}$. 

To prove that $Z_{\infty,k,\infty}$ converges as $k\to\infty$, we show that it is a martingale:
\begin{align*}
    \E\left[Z_{\infty,k+1,\infty}\ \middle|\ \left(\Delta(\bm{i})\right)_{\bm{i}\in \bigcup_{1 \le\ell\le k}\mathbb{N}^\ell}\right]
    &=\sum_{\bm{i}\in \mathbb{N}^k}\exp\left(-\sum_{\bm{j}\in P(\bm{i})}\Delta(\bm{j})\right)\cdot\E\left[\sum_{a=1}^\infty\exp\left(-\sum_{1 \le b\le a}\Delta(\bm{i},b)\right)\right]\\
    &=\sum_{\bm{i}\in \mathbb{N}^k}\exp\left(-\sum_{\bm{j}\in P(\bm{i})}\Delta(\bm{j})\right)\cdot\sum_{a = 1}^{\infty}2^{-a}\\
    &=\sum_{\bm{i}\in \mathbb{N}^k}\exp\left(-\sum_{\bm{j}\in P(\bm{i})}\Delta(\bm{j})\right)\cdot1=Z_{\infty,k,\infty}.
\end{align*}
In the second line, we used that if $\Delta \sim \text{Exp}(1),$ then $\E[\exp(-\Delta)] = \frac{1}{2}$, and thus by independence $\E\left[\exp\left(-\sum_{b=1}^{a}\Delta(\bm{i}, b) \right)\right] = 2^{-a}$. Since $Z_{\infty,k,\infty}$ is a nonnegative martingale, Doob's martingale theorem ensures that it converges almost surely to some random variable $Z$.

Next, we show that $Z\sim\text{Exp}(1)$.
To do so, we first rewrite $Z_{\infty,k,\infty}$. Note that $1\in P(\bm{i})$ for all $\bm{i}\in\mathbb{N}^k$. Therefore,
\begin{align*}
Z_{\infty,k,\infty}
&=e^{-\Delta(1)}\sum_{\bm{i}\in\N^k}\exp\left(-\sum_{\bm{j}\in P(\bm{i})\setminus\{1\}}\Delta(\bm{j})\right)\\
&=e^{-\Delta(1)}\sum_{\bm{i}\in\N^k\ :\ i_1=1}\exp\left(-\sum_{\bm{j}\in P(\bm{i})\setminus\{1\}}\Delta(\bm{j})\right)+e^{-\Delta(1)}\sum_{\bm{i}\in\N^k\ :\ i_1>1}\exp\left(-\sum_{\bm{j}\in P(\bm{i})\setminus\{1\}}\Delta(\bm{j})\right)\\
&\stackrel{D}=e^{-\Delta(1)}Z_{\infty,k-1,\infty}+e^{-\Delta(1)}Z_{\infty,k,\infty}.
\end{align*}
Since $\Delta(1)\sim\text{Exp}(1)$, we have $e^{-\Delta(1)}\sim\text{Unif}([0,1])$. Taking the limit $k\to\infty$ on both sides tells us that
\[
Z\stackrel{D}=U\cdot(Z_{1}'+Z_{2}'),
\]
where $U\sim\text{Unif}([0,1])$ and $Z_{1}',Z_{2}'$ are independent copies of $Z$. We now derive the moment-generating function (MGF) of $Z$.
\begin{align*}
    M(t)=\E[e^{-tZ}]=\E\left[\E\left[e^{tU\cdot(Z_1'+Z_2')}\ \middle|\ Z_1',Z_2'\right]\right]=\E\left[\int_0^1e^{tu\cdot(Z_1'+Z_2')}du\right]=\E\left[\frac{e^{t\cdot(Z_1'+Z_2')}-1}{t\cdot(Z_1'+Z_2')}\right].
\end{align*}
Taking the derivative with respect to $t$ yields
\begin{align*}
    M'(t)
    =\E\left[\frac{(Z_1'+Z_2')\cdot e^{t\cdot(Z_1'+Z_2')}}{t\cdot(Z_1'+Z_2')}\right]-\E\left[\frac{e^{t\cdot(Z_1'+Z_2')}-1}{t^2\cdot(Z_1'+Z_2')}\right]=\frac{M(t)^2-M(t)}{t}.
\end{align*}
This is a separable differential equation:
$$
\frac{M'(t)}{M(t)^2-M(t)}=\frac1t.
$$
Integrating both sides yields
$$
\log(1-M(t))-\log M(t)=\log(t)+c\\
\Rightarrow \frac{1-M(t)}{M(t)}=te^c\\
\Rightarrow M(t)=\frac{1}{1+te^c}.
$$
The constant $e^c$ is determined by $1=\mathbb{E}[Z]=M'(0)=-e^{c}$, so that $e^c = -1$. We conclude that
\(
\E[e^{tZ}]=(1-t)^{-1},
\)
so that indeed $Z\sim\text{Exp}(1)$.

Finally, we prove $Z_{\infty,k,r_k}\stackrel\P\to Z$ for our choice of $r_k$.
Since $Z_{\infty,k,r}$ is increasing in $r$, we have $Z_{\infty,k,\infty}-Z_{\infty,k,r_k}>0$.
The Markov inequality allows us to bound
\[
\P(Z_{\infty,k,\infty}-Z_{\infty,k,r_k}>\varepsilon)\le\frac1\varepsilon\left(1-(1-2^{-r_k})^k\right)=\bigO(k2^{-r_k})\to0.
\]
This proves $Z_{\infty,k,\infty}-Z_{\infty,k,r_k}\stackrel\P\to0$, so that $Z_{\infty,k,r_k}\stackrel\P\to Z$.
\end{proof}

\end{document}
